\documentclass[11pt]{article}
\usepackage{preamble}
\setlength{\parskip}{4pt}

\begin{document}

\daybanner{AP Statistics \textperiodcentered\ Topic 1.9 (CED 1.9) \textperiodcentered\ B: 2026-09-24 \textperiodcentered\ E: 2026-09-30 \textperiodcentered\ TEACHER KEY}{Which Class Did Better?}

\sectionbanner{CED TETHER}

{\small
\begin{itemize}[leftmargin=1.5em,itemsep=2pt,topsep=2pt]
  \item \textbf{Skill 2.D} --- Compare distributions or relative positions of points within a distribution.
  \item \textbf{Enduring Understanding UNC-1} --- Graphical representations and statistics allow us to identify and represent key features of data.
  \item \textbf{LO UNC-1.N} --- Compare graphical representations for multiple sets of quantitative data. [Skill 2.D]
  \item \textbf{EK UNC-1.N.1} --- Any of the graphical representations, e.g., histograms, side-by-side boxplots, etc., can be used to compare two or more independent samples on center, variability, clusters, gaps, outliers, and other features.
  \item \textbf{LO UNC-1.O} --- Compare summary statistics for multiple sets of quantitative data. [Skill 2.D]
  \item \textbf{EK UNC-1.O.1} --- Any of the numerical summaries (e.g., mean, standard deviation, relative frequency, etc.) can be used to compare two or more independent samples.
\end{itemize}
}
\textbf{Essential question:} What does `did better' mean for a whole class, and which feature of a plot answers it?\par
\textbf{DOK-3 skill:} 4.B \textperiodcentered\ \textbf{FRQ pattern:} \texttt{compare-two-boxplots-adjudicate-competing-claims-then-choose-one-statistic}\par
\textbf{DOK rationale:} Part (c) requires weighing two competing claims about which class did better against several features of two boxplots at once, then committing to a single statistic and defending it; either claim can be supported by one feature alone, so the answer depends on coordinating center, variability, and an outlier rather than on any single procedure.\par

\frameworkphaseheader{Do Now (first take) $\rightarrow$ Explore (video + follow-along) $\rightarrow$ Exit (finish + turn in)}{1 $\rightarrow$ 3}{45}{%
\begin{itemize}[leftmargin=1.5em,itemsep=2pt,topsep=2pt]
  \item Board slide up at the bell. Read Sam's and Tia's lines aloud once; take a quick hand-count of first takes but do not react.
  \item Video follow-along from minute 5.
  \item Board slide back up at minute 33. Circulate with the prompts; insist on comparative sentences.
  \item Collect at the door. Score part (c) only, E/P/I.
\end{itemize}
}{%
\begin{itemize}[leftmargin=1.5em,itemsep=2pt,topsep=2pt]
  \item One-sentence first take naming one plot feature.
  \item Video follow-along (comparing distributions: center, variability, shape, unusual features).
  \item Finish (a)--(c); exit line; turn in.
\end{itemize}
}{%
\begin{itemize}[leftmargin=1.5em,itemsep=2pt,topsep=2pt]
  \item How many students does the top of a whisker represent? How many does the median line represent?
  \item If the 40 in Period B were a 33 instead, would Sam's claim change? Would Tia's?
  \item Point at the box widths. Which class would you rather be a random student in? Why?
  \item Your parent asks `how did the class do?' --- which number answers that question?
\end{itemize}
}{Solo teacher. Circulate ~5 pairs. Answer questions with questions; do not confirm Tia until two features are cited.}

\begin{callout}{\IconWarn\ WATCH FOR}{calloutred}
\begin{itemize}[leftmargin=1.5em,itemsep=2pt,topsep=2pt]
  \item Sentences that describe one plot at a time with no comparison word.
  \item Reading the whisker end as `most students'.
  \item Choosing the mean for Period B without noticing the outlier at 5.
\end{itemize}
\end{callout}

\sectionbanner{SCORING PART (c) --- E / P / I}

\textbf{Expected elements}
\begin{itemize}[leftmargin=1.5em,itemsep=2pt,topsep=2pt]
  \item \textbf{picks-tia} (required) --- Chooses Tia's claim as better supported
  \item \textbf{cites-median-comparison} (required) --- Cites the medians (31 vs. 27) or equivalent center comparison as a settling feature
  \item \textbf{cites-spread-or-outlier} (required) --- Cites the IQR / range / outlier as a second settling feature
  \item \textbf{median-over-max-reason} (required) --- Chooses the median (or another center measure) and explains that the maximum describes one student, not the class
\end{itemize}

{\small\renewcommand{\arraystretch}{1.25}
\noindent\begin{tabular}{|p{0.5in}|p{5.9in}|}
\hline
\textbf{E} & Picks Tia, cites center AND variability (or the outlier), chooses the median with the `one student vs. the class' reasoning stated in context. \\ \hline
\textbf{P} & Picks Tia with only one feature, OR chooses the median without contrasting it with the maximum, OR gives the reasoning but names the mean without addressing the outlier. \\ \hline
\textbf{I} & Picks Sam, or compares with no specific plot features, or chooses the maximum. \\ \hline
\end{tabular}}

\textbf{Common mistakes}
\begin{itemize}[leftmargin=1.5em,itemsep=2pt,topsep=2pt]
  \item Reading the maximum as `the class did better'
  \item Describing each boxplot separately instead of comparing (no comparative language)
  \item Calling the wider box `more scores' rather than `more variable'
  \item Choosing the mean for Period B without noticing the outlier at 5 pulls it down
\end{itemize}

\bankitem{Optional \#1 --- not collected}{\dokbadge{2}~Two bus routes were timed for 20 mornings each. Route 1: mean 24.0 min, SD 6.5 min, median 22 min. Route 2: mean 24.5 min, SD 1.8 min, median 24 min. Write two sentences comparing the routes for a student who must not be late, using comparative language and the statistics given.}{}
\answer{Centers are nearly the same (medians 22 vs 24, means 24.0 vs 24.5). Route 2 is far more consistent (SD 1.8 vs 6.5 min); Route 1's mean above its median suggests some very long mornings. A student who must not be late should prefer Route 2: a typical morning is about as long, but the bad mornings are much shorter.}

\clearpage
\focusbanner{TODAY'S PROBLEM --- annotated}

Two AP Statistics classes took the same 40-point quiz. The side-by-side boxplots show the scores. Two students read them differently. \textbf{Sam:} ``Period B did better --- its highest score is 40 and Period E's is only 38.'' \textbf{Tia:} ``Period E did better --- its median is higher and its scores are more consistent.''\par
\begin{center}
\begin{tikzpicture}
\begin{axis}[
  width=0.68\linewidth,
  height=1.60in,
  ytick={2,1}, yticklabels={{Period B},{Period E}}, ymin=0.4, ymax=2.6,
  axis x line=bottom, axis y line*=left, y axis line style={draw=none}, boxplot/box extend=0.5, xmajorgrids, enlarge x limits=0.06,
  xlabel={Quiz score (out of 40)}, tick label style={font=\small}, label style={font=\small},
  ytick style={draw=none}, yticklabel style={font=\small\bfseries}
]
\addplot[boxplot prepared={lower whisker=12, lower quartile=22, median=27, upper quartile=33, upper whisker=40, draw position=2}, fill=calloutblue, draw=dayblue] coordinates {};
\addplot[only marks, mark=*, color=warmred] coordinates {(5,2)};
\addplot[boxplot prepared={lower whisker=18, lower quartile=26, median=31, upper quartile=34, upper whisker=38, draw position=1}, fill=calloutblue, draw=dayblue] coordinates {};
\end{axis}
\end{tikzpicture}
\end{center}
\begin{firsttakebox}{FIRST TAKE (5 min)}
Which class did better on the quiz? One sentence, and name ONE feature of the plot you used.\par
\answer{Not graded. Expect a split toward Sam (the 40 is eye-catching); that split is the lesson.}
\end{firsttakebox}

\textit{Rules callout ``COMPARING DISTRIBUTIONS (keep on the page)'' is printed on the student sheet.}\par\medskip

\dokbadge{1}~\textbf{(a)} Read off the median and the maximum score for each class.\par
\answer{Period B: median 27, maximum 40. Period E: median 31, maximum 38.}

\dokbadge{2}~\textbf{(b)} Compare the two distributions of quiz scores in context. Use comparative language (higher / lower, more / less variable) and address center, variability, shape, and any unusual features.\par
\answer{Center: Period E's median (31) is higher than Period B's (27). Variability: Period E is less variable --- IQR 34 - 26 = 8 vs. Period B's 33 - 22 = 11; Period B's range is also wider (12 to 40 vs. 18 to 38), and Period B has a low outlier at 5. Shape: Period B is roughly symmetric within the box; Period E's box is slightly skewed left (the median sits nearer Q3). Unusual: the single low score of 5 in Period B. In context: a typical Period E student scored a few points higher and the class was more consistent, while Period B had both the highest and the lowest scores.}

\dokbadge{3}~\textbf{(c)} Whose claim is better supported, Sam's or Tia's? Cite two specific features of the boxplots that settle it. Then: if you could report only ONE statistic per class to a parent asking `how did the class do?', which statistic would you choose, and why that one rather than the maximum?\par
\answer{Tia's. Two settling features: the medians (31 vs. 27 --- the typical student in E scored higher) and the spread (E's IQR 8 vs. B's 11, plus B's outlier at 5). Sam's evidence is a single value: one student in B scored 40, which says nothing about the class. Report the median: it describes the typical student and is not moved by one very high or very low score; the maximum describes exactly one student and would make Period B look better because of one paper.}

\begin{summaryexitbox}{EXIT}
One thing the video changed about my first take: \hrulefill
\end{summaryexitbox}

\end{document}
